In this section we give the basic concepts and definitions of coding theory (\cite{macwilliams1977theory}), required to formally define
the problems on which HQC and similar cryptosystems are based.

\subsection{General definitions}
We denote with $\mathds{F}_2$ the binary finite field and with $\mathds{Z}$ the integer ring.\\
We define $\mathcal{V} = \mathds{F}_2^n$ the vector space of dimension $n$ over $\mathds{F}_2$ 
and $\mathcal{R}=\mathds{F}_2[X]/(X^n-1)$ the ring of binary polynomials of degree less than $n$ for some $n \in \mathds{N}$.\\
Additionally, we denote by $\texttt{wt}(\cdot)$ the Hamming weight of a vector \textit{i.e.} the number of its non-zero coordinates.

\begin{definition}[Circulant matrix]
    Let $\mathbf{v} = (v_0, \cdots, v_{n-1}) \in \mathcal{V}$. The circulant matrix induced by $\mathbf{v}$ is defined as
    a matrix whose $i$-th row is the vector $\mathbf{v}$ shifted by $i$ positions to the right:
    \begin{equation*}
        \mathbf{circ(v)} = \begin{pmatrix}
            v_0 & v_1 & \cdots & v_{n-1} \\
            v_{n-1} & v_0 & \cdots & v_{n-2} \\
            \vdots & \vdots & \ddots & \vdots \\
            v_1 & v_2 & \cdots & v_0
            \end{pmatrix} \in \mathds{F}_2^{n\times n}
    \end{equation*}
\end{definition}
The algebra of $n\times n$ circulant matrices over $\mathds{F}_2$ is isomorphic to the algebra of polynomials
in the ring $\mathcal{R}$:
\begin{align*}
    \mathbf{circ(v)} \simeq v_0 + v_1x + v_2x^2 + \cdots + x_{n-1}x^{n-1} \\
\end{align*}

\begin{definition}[Linear $q$-ary Codes]
    A Linear $q$-ary code $\mathcal{C}$ of lenght $n$ and dimension $k$ (denoted as $[n, k]$) is a subspace of $\mathds{F}_q^n$.\\
    Due to their vast use in cryptography (and in particular in this work), the following definitions and results are relative
    to binary codes ($q=2$).
\end{definition}

\begin{definition}[Linear Codes]
    A linear code $\mathcal{C}$ of lenght $n$ and dimension $k$ (denoted as $[n, k]$) is a subspace of $\mathcal{V}$ 
    of dimension $k$; the elements of $\mathcal{C}$ are called \textit{codewords}.\\
    Due to the linearity of $\mathcal{V}$, any linear combination of two codewords is also a codeword.
\end{definition}

\begin{definition}[Generator Matrix]
    Let $\mathcal{C} = [n,k]$ be a linear code. The matrix $\mathbf{G} \in \mathds{F}_2^{k\times n}$ is called the generator 
    matrix of $\mathcal{C}$ if
    \begin{equation*}
        \mathcal{C} = \{\mathbf{mG}, \text{ for } \mathbf{m}\in \mathds{F}_2^k\}
    \end{equation*}
\end{definition}

\begin{definition}[Parity-Check Matrix]
    Let $\mathcal{C} = [n,k]$ be a linear code. The matrix $\mathbf{H} \in \mathds{F}_2^{(n-k)\times n}$ is called the 
    parity-check matrix of $\mathcal{C}$ if
    \begin{equation*}
        \mathcal{C} = \{\mathbf{v} \in \mathds{F}_2^n \text{ such that } \mathbf{Hv}^\top = 0\}
    \end{equation*}
    In other words, $\mathbf{H}$ is the generator matrix of the dual code $\mathcal{C}^\perp$.
\end{definition}

\begin{definition}[Rank]
    Let $\mathcal{C} = [n,k]$ be a linear code. The rank of $\mathcal{C}$ is defined as the dimension of the space spanned
    by its codewords or, in other words, the rank of its parity-check matrix.
    \begin{equation*}
        r(\mathcal{C}) = r(\mathbf{H}) = n-k
    \end{equation*}
\end{definition}
A high rank code can correct more errors with fewer parity-check equations, but it usually is harder to construct and decode.

\begin{definition}[Syndrome]
    Let $\mathcal{C} = [n,k]$ be a code, $\mathbf{H}$ its parity matrix and $\mathbf{v} \in \mathds{F}_2^n$ a word. 
    The syndrome of $\mathbf{v}$ is defined as
    \begin{equation*}
        \sigma(\mathbf{v}) = \mathbf{Hv}^\top
    \end{equation*}
    From this definition, we can say that a vector $\mathbf{v}$ is a codeword if and only if its syndrome is the zero vector.
\end{definition}

\begin{definition}[Rate]
    Let $\mathcal{C} = [n,k]$ be a code. The rate of $\mathcal{C}$ is defined as $R = \frac{k}{n}$.
\end{definition}
The rate of a code is a measure of its efficiency: with a higher rate we can transmit more information with the same amount
of bits, at the cost of reduced redundancy (and thus reduced error correcting capability).

\begin{definition}[Minimum distance]
    Let $\mathcal{C} = [n,k]$ be a code and $\omega$ a norm on $\mathcal{V}$. The minimum distance of $\mathcal{C}$ is defined as
    \begin{equation*}
        d = \min_{\mathbf{v}, \mathbf{w} \in \mathcal{C}, \mathbf{v}\neq \mathbf{w}} \omega(\mathbf{v} - \mathbf{w})
    \end{equation*}
\end{definition}
The length, dimension and minimum distance of a linear code $\mathcal{C} = [n,k,d]$ are tied by the Singleton bound~\cite{singleton1964maximum}:
\begin{equation*}
    d \leq n-k+1
\end{equation*}
If the equality holds, $\mathcal{C}$ is called a maximum distance separable (MDS) code, and it has desirable properties in terms of error correction.

\begin{definition}[Error Correction Capability]
A linear code $\mathcal{C} = [n,k,d]$ can decode arbitrary codewords with at most $\delta = \lfloor \frac{d-1}{2}\rfloor$ errors.
If $\mathcal{C}$ is a MDS code, we can rewrite this as $\delta = \lfloor \frac{n-k}{2} \rfloor$.
\end{definition}

\begin{definition}[Gilbert-Varshamov bound]
    Let $\mathcal{C} = [n,k]$ be a linear code. The Gilbert-Varshamov bound~\cite{gilbert1952comparison,varshamov1957estimate} 
    states that $\mathcal{C}$ can correct at least $t$ errors if its size satisfies
    \begin{equation}
        2^k \geq \sum_{i=0}^t \binom{n}{i} \label{eq:gv_bound}
    \end{equation}
\end{definition}

\subsection{Code families}
Informally, when we say that a code is \textit{good}, we mean that it satisfies some properties, such as:
\begin{itemize}
    \item High minimum distance: the higher the minimum distance, the more errors the code can correct.
    \item High rate: the higher the rate, the more information we can transmit with the same amount of bits.
    \item Efficiency: we want the encoder and decoder to be fast and simple to implement.
\end{itemize}
In this section we will describe some of the most common code families and the properties they satisfy.

\begin{definition}[Cyclic Codes]
    Let $\mathbf{c} = (c_0, \cdots, c_{n-1}) \in \mathcal{V}$ be a vector. A linear code $\mathcal{C} = [n, k]$ is Cyclic
    if, for any $\mathbf{v} \in \mathcal{C}$, the vector obtained by applying a circular shift to $\mathbf{v}$ is also a codeword.
    codeword. If we think of $\mathbf{v}$ as a polynomial in $\mathcal{R}$, we can say that $\mathcal{C}$ is cyclic if,
    for any $\mathbf{v} \in \mathcal{C}$, the polynomial $X \cdot \mathbf{v}$ is also a codeword.
\end{definition}

\begin{definition}[Quasi-Cyclic Codes]
    Let $\mathbf{c} = (\mathbf{c}_0, \cdots, \mathbf{c}_{s-1}) \in \mathds{F}_2^{sn}$ be a vector, where we consider 
    each $\mathbf{c}_i$ as a vector in $\mathds{F}_2^n$.
    A linear code $\mathcal{C} = [sn,k]$ is Quasi-Cyclic of index $s$ if, for any $\mathbf{c} \in \mathcal{C}$, 
    the vector obtained by applying a circular shift to every block $\mathbf{c}_i$ simultaneously is also a codeword.
    QC codes are a generalization of cyclic codes.
\end{definition}

\begin{definition}[Double Circulant Codes]
    A code $\mathcal{C}$ is double circulant if it is Quasi-Cyclic of even length and index $s=2$.
\end{definition}

\subsubsection{Random codes}
    In a linear random code the codewords are chosen at random; specifically, we can build the parity matrix $\mathbf{H}$ of
    the code by randomly sampling each entry from a Bernoulli distribution.\\
    For sufficiently large $n$, random linear codes reach the Gilbert-Varshamov bound; moreover, the probability of decoding
    error decreases exponentially as the minimum distance increases. On the other hand, the computational complexity of a decoder
    for a random code grows exponentially with the length of the code.

\subsubsection{Repetition codes}
    Repetition codes are among the simplest codes, as they are just a repetition of the same vector $k$ times.
    Repetition codes can be easily decoded by majority voting, making the decoding process very fast. However,
    they have a very low rate and they are not efficient in terms of error correction capability.

\subsubsection{Concatenated codes}
    Concatenated codes~\cite{forney1965concatenated} are built by combining an \textit{inner} code $\mathcal{C}_{in}$ and an
    \textit{outer} code $\mathcal{C}_{out}$. The goal of concatenated codes is to obtain a code which has both a 
    polynomial-time decoder and exponentially decreasing error probability with increased block length.

\subsubsection{Reed-Muller codes}
    Reed-Muller codes~\cite{reed1953class,muller1954application} are a family of codes; they are usually denoted as \textsf{RM}$(r,m)$, 
    as they have a message length which is equal to
    \begin{equation*}
        k = \sum_{i=0}^r \binom{m}{i}
    \end{equation*}
    RM codes are:
    \begin{itemize}
        \item Locally testable: it is possible to determine whether a string is a word in the code by testing only a few bits of the string.
        \item Locally decodable: it is possible to recover the original information by using a small amount of bits of the codeword.
        \item List decodable: TODO
    \end{itemize}

    Due to their properties, we can describe RM codes based on low degree polynomials.\\







\subsection{Difficult problems for cryptography}
The problems we are interested in are all variants of the \textit{decoding problem}, which requires to find 
the closest codeword to a given word with respect to a certain metric (the Hamming weight in our case).\\
In the case of linear codes we have shown that the decoding problem stays the same whether one is given 
the syndrome of the received vector or the vector itself; for this reason, we will only consider 
the \textit{Syndrome Decoding} (\textsf{SD}) problem.\\

\begin{definition}[\textsf{SD} Distribution]
    Given $n,k,w \in \mathds{N}$, the distribution $\textsf{SD}(n, k, w)$ randomly samples a parity-check matrix 
    $\mathbf{H} \xleftarrow{\mathdollar} \mathds{F}_2^{(n-k)\times n}$, a vector 
    $\mathbf{x} \xleftarrow{\mathdollar} \mathds{F}_2^n$ such that $\texttt{wt}(\mathbf{x}) = w$ 
    and returns ($\mathbf{H}, \sigma(\mathbf{x})$).
\end{definition}

\begin{definition}[\textsf{SD} Problem]
    Given $(\mathbf{H}, \mathbf{y}^\top) \in \mathds{F}_2^{(n-k)\times n} \times mathds{F}_2^{(n-k)}$ from the $\textsf{SD}(n, k, w)$ 
    distribution, the \textsf{SD} problem requires to find a vector $\mathbf{x} \in \mathds{F}_2^n$ such that 
    $\sigma(\mathbf{x}) = \mathbf{y}^\top$ and $\texttt{wt}(\mathbf{x}) = w$.
\end{definition}
This problem has been proven to be NP-complete in~\cite{berlekamp1978inherent}.

\begin{definition}[Decisional \textsf{SD} Problem]
    Given $(\mathbf{H}, \mathbf{y}^\top) \in \mathds{F}_2^{(n-k)\times n} \times mathds{F}_2^{(n-k)}$, the \textsf{DSD} problem 
    requires to determine whether $(\mathbf{H}, \mathbf{y}^\top)$ comes from the \textsf{SD} distribution or from the random
    distribution over $\mathds{F}_2^{(n-k)\times n} \times mathds{F}_2^{(n-k)}$.
\end{definition}


 



